\section{The negative phase and actual point labels}\label{sec:phase}

The sign $p$ in \eqref{eq:sigma} changes with the parity of $R$.
It changes the central cycle structure, but not the generic bulk
multiplicities. To keep the ordinary clock and the point coordinates
correct, introduce
\[
 S(x,y)=(-x,-y),\qquad C(x,y)=(-x,y).
\]

\begin{proposition}[Phase transfer]\label{prop:phase}
The identities
\begin{equation}\label{eq:phase}
 Sg=gS,\qquad g_-=C(Sg)C
\end{equation}
hold on $\Z^2$. Let $\mathcal K_+$ be the set described by the
clipped cells of Proposition~\ref{prop:core}. The set of points whose
entire $g_-$ orbit lies in $B_0$ is $C\mathcal K_+$. For
$z\in\mathcal K_+$, its actual negative-phase orbit is
\begin{equation}\label{eq:point-conversion}
 g_-^n(Cz)=C S^n g^n z.
\end{equation}
All generic periods $4,12,20$ and their multiplicities are unchanged.
The positive central cycles $1:1,5:2,6:1$ become
$1:1,3:2,10:1$ in the negative phase.
\end{proposition}
\begin{proof}
Oddness of $\sine$ proves $Sg=gS$; direct substitution proves the
second identity. Both $S$ and $C$ preserve $B_0$. Since
$(Sg)^n=S^ng^n$, a point has its entire $Sg$ orbit in $B_0$
if and only if it has its entire $g$ orbit there. Conjugating by
$C$ proves the set statement and \eqref{eq:point-conversion}.

Suppose that $z$ has $g$-period $\ell\in\{4,12,20\}$. Because
$\ell$ is even, $(Sg)^\ell z=z$. If $(Sg)^nz=z$ with $n$
even, then $g^nz=z$, so $\ell\mid n$. If $n$ is odd, then
$g^nz=-z$, which implies $g^{2n}z=z$ and hence $\ell\mid2n$.
This is impossible because $4\mid\ell$ but $n$ is odd.
Thus its least $Sg$ period is exactly $\ell$. The number of points
of each generic period, and consequently the cycle multiplicity,
does not change.

Negation exchanges the two cycles in \eqref{eq:five-cycle}, so $Sg$
joins their ten points into one $10$-cycle. On the cycle
\eqref{eq:six-cycle}, negation equals $g^3$; therefore $Sg=g^4$
there, producing two $3$-cycles. The origin is still fixed.
Applying $C$ to these orbit labels yields the claimed negative-phase
cycles for $g_-$.
\end{proof}

For example, a negative-phase central $10$-cycle is obtained by
applying $C S^n$ to the $n$th point of \eqref{eq:five-cycle}, with
the latter point read modulo $5$, for $0\le n<10$. This is an
explicit reconstruction of ten actual points. Directly reusing
the positive-phase labels without $C$ would not be the stated
conjugacy. Because $p=+1$ precisely when $R$ is odd,
Proposition~\ref{prop:phase} proves Table~\ref{tab:central} and the
bulk part of Theorem~\ref{thm:main} for the actual map $h_d$.

